\section{Method}
\label{sec:method}

IAM is a single causal language model, fully fine-tuned from
Qwen3.5-2B \citep{qwen3}, that (a) renders an agent trajectory into a
canonical token stream and (b) emits a verdict after every decision step. We
first describe the canonical schema, then the dense rendering and training
objective, and finally the online \texttt{check()} semantics used at
deployment.

\subsection{Canonical trajectory schema}
\label{sec:method:schema}

We represent every agent run, regardless of framework, as a canonical
\texttt{Trajectory}: a task instruction plus an ordered list of
\texttt{TraceStep}s. Each step carries a \texttt{role}
($\{$user, agent, tool, system$\}$), an optional \texttt{thought}, an optional
structured \texttt{action} (e.g. \texttt{\{"tool":"Read","args":\{"path":".env"\}\}}),
and an optional \texttt{observation}. A \emph{decision step} is any agent step
--- the point at which the monitor runs in deployment.

A step renders to a flat text block using bare special tokens:
\begin{quote}\small\ttfamily
<step> idx=$k$ role=agent\\
<thought>...</thought>\\
<action>Read(\{'path': '.env'\})</action>\\
<observation>...</observation>\\
</step>
\end{quote}
Per-step metadata (index, role, tool) is emitted as text \emph{after} the
opening tag rather than as XML attributes, so that the structural tags
(\texttt{<step>}, \texttt{<action>}, \dots) are matched exactly as added special
tokens by the tokenizer. The full taxonomy carried by a trajectory's label
(\texttt{risk\_source} $\times$ \texttt{failure\_mode} $\times$
\texttt{harm\_category}) is taken directly from ATBench \citep{atbench}
(Section~\ref{sec:data}).

\subsection{Dense per-step verdict rendering}
\label{sec:method:dense}

The key training idea is \emph{dense per-step verdict supervision}: after every
decision step we emit a verdict slot,
\begin{quote}\small\ttfamily
<verdict><symbol>OK</symbol></verdict>
\end{quote}
for safe steps, and at the anomaly step the full triple
\begin{quote}\small\ttfamily
<verdict><symbol>STOP</symbol><type>FAILURE\_MODE</type><reason>...</reason></verdict>
\end{quote}
Safe steps are supervised with the symbol only; the anomaly step additionally
supervises a risk \texttt{<type>} (the ATBench failure mode) and a
natural-language \texttt{<reason>}. Crucially, the per-step verdicts are derived
automatically from the single trace-level anomaly label --- no manual per-step
annotation is required. When the labelled \texttt{anomaly\_step} is not itself a
decision step, the \textsc{STOP} is attached to the last decision step at or
before it; when the anomaly step is unspecified (e.g. some content-harm traces),
\textsc{STOP} is attached to the final decision step so that the trace still
carries supervision to flag by its end.

Supervising every decision step --- not just the anomaly --- means the model
sees many \textsc{OK} decisions in attack and benign traces alike, so it learns
the negative case (``do not stop here'') as strongly as the positive one. This
directly targets over-flagging, which is the failure mode that drives the
false-stop rate in deployment.

\paragraph{Truncation at the anomaly.}
Because \textsc{STOP} is terminal --- in deployment, execution halts at the
flagged action --- training trajectories are truncated \emph{after} the anomaly
step. The model is never shown post-anomaly steps, matching the deployment
regime where it never sees what would have happened after a block.

\subsection{Training objective}
\label{sec:method:loss}

We fine-tune with two masked next-token objectives over the rendered stream:
(i) a self-supervised next-action term over tokens inside \texttt{<action>}
spans (learning the distribution of normal agent behavior), and (ii) a
supervised verdict term over tokens inside \texttt{<verdict>} spans (the dense
\textsc{OK}/\textsc{STOP} decisions, plus type/reason at the anomaly). All
tokens outside these spans are masked out of the loss ($-100$ labels).

\paragraph{STOP up-weighting.}
The dense layout is highly imbalanced: a trajectory with $N$ decision steps
contributes $N{-}1$ \textsc{OK} verdicts and at most one \textsc{STOP}, giving an
$\textsc{OK}{:}\textsc{STOP}$ ratio of roughly $7{:}1$ in our data. Left
uncorrected, the model converges to ``never stop.'' We therefore up-weight the
cross-entropy on the \textsc{STOP} symbol token by a factor $W_{\textsc{STOP}}=8$,
chosen to approximately balance the two classes. This is the single most
important hyperparameter for moving the prevention rate off the floor.

\subsection{Left-truncation for a bounded recent window}
\label{sec:method:lefttrunc}

Trajectories can exceed the model's context. We cap sequences at
\texttt{max\_seq\_len}$=4096$ tokens and, critically, \emph{left}-truncate
(keep the tail) rather than the default right-truncation. The terminal
\textsc{STOP} verdict lives at the end of the dense stream, so right-truncation
would drop exactly the label we need; keeping the tail preserves the anomaly and
its verdict, and matches a deployment monitor's bounded recent-context window.
A partial span left dangling at the new start is simply left unsupervised (its
span depth never opens), so truncation never corrupts the mask.

\subsection{Online inference: \texttt{check()}}
\label{sec:method:check}

At deployment the monitor sits in the agent's execution loop. Before dispatching
a proposed action, the executor calls:
\begin{quote}\small\ttfamily
verdict = monitor.check(trace\_so\_far, next\_action)\\
if verdict.block:  \# verdict.symbol == "STOP"\\
\hphantom{xx}raise SecurityError(verdict.risk\_type, verdict.explanation)
\end{quote}
\texttt{check()} appends the proposed action as a hypothetical agent step,
renders the trajectory in the dense layout with an \textsc{OK} verdict after
every \emph{prior} decision step (consistent with the fact that, had any earlier
step been \textsc{STOP}, execution would already have halted), and leaves the
final verdict open --- the string ends at \texttt{<verdict><symbol>}. The model
runs a single forward pass; we take the next-token logits, restrict to the first
token of each of $\{\textsc{OK},\textsc{WARN},\textsc{STOP}\}$, and softmax to
obtain $P(\textsc{STOP})$. The action is blocked iff $P(\textsc{STOP})$ exceeds a
calibrated threshold $\tau$. Only when the verdict is \textsc{STOP} do we run a
short greedy decode to produce the risk \texttt{<type>} and \texttt{<reason>}, so
the common (\textsc{OK}) path costs a single forward pass. The same left-tail
truncation as in training is applied so the open verdict always survives and the
scored $P(\textsc{STOP})$ stays in-distribution.

\paragraph{Operating-point calibration.}
$\tau$ is a deployment knob, not a fixed constant. We calibrate it on held-out
data against a target false-stop rate (Section~\ref{sec:experiments}); both the
prevention rate and the false-stop rate are monotonically non-increasing in
$\tau$, so the smallest $\tau$ meeting the FSR constraint also maximizes
prevention.

\paragraph{Baselines reuse the same harness.}
For comparison we retain a legacy action-span perplexity scorer (next-token PPL
restricted to the last \texttt{<action>} span) and a keyword/regex scorer; both
emit one score per step under the same interface, so they slot into the same
evaluation harness as the verdict head (Section~\ref{sec:experiments}).
